\documentclass[11pt]{article}
% Standalone preamble for web-only supplementary notes.
\usepackage[T1]{fontenc}
\usepackage{amsmath,amssymb,amsthm}
\usepackage{hyperref}
\usepackage{cleveref}
\DeclareMathOperator{\Map}{Map}
\DeclareMathOperator{\Fun}{Fun}
\newcommand{\iso}{\xrightarrow{\sim}}
\newtheorem{lemma}{Supplementary Lemma}
\newtheorem{corollary}[lemma]{Supplementary Corollary}
\crefname{lemma}{Supplementary Lemma}{Supplementary Lemmas}
\Crefname{lemma}{Supplementary Lemma}{Supplementary Lemmas}
\crefname{corollary}{Supplementary Corollary}{Supplementary Corollaries}
\Crefname{corollary}{Supplementary Corollary}{Supplementary Corollaries}

\begin{document}
\section*{Substitution in morphism expressions}

The \href{segal-axiom.html}{Segal section} uses composition in a common
parameter category. The supporting calculation is that restriction and
changes of endpoints preserve the chosen composite, with its endpoint data.

Write $\operatorname{Ar}(C)=\Fun([1],C)$, with endpoint evaluations $e_0,e_1$.
A morphism expression from $x$ to $y$, for $x,y\colon\Gamma\to C$,
consists of $A\colon\Gamma\to\operatorname{Ar}(C)$ and frames
$p\colon e_0A=x$, $q\colon e_1A=y$.
An endpoint-preserving comparison from $(A,p,q)$ to $(A',p',q')$ is an
identification $\eta\colon A=A'$ together with
\[
 p'\cdot(e_0\star\eta)=p,\qquad q'\cdot(e_1\star\eta)=q.
\]
Retargeting along $a\colon x=x'$ and $b\colon y=y'$ replaces the frames
$p,q$ by $a\cdot p,b\cdot q$ and leaves the arrow family unchanged.
We use $\star$ for whiskering and $\cdot$ for vertical composition,
rightmost factor first. All categories involved are fixed absolute categories.

\begin{lemma}
\label[lemma]{lem:Supplementary_Expression_Restriction}
For $r\colon\Delta\to\Gamma$, restriction takes endpoint-preserving
comparisons to endpoint-preserving comparisons. The restricted arrow is
$Ar$ and its source frame is
\[
 (p\star r)\cdot\alpha_{r,A,e_0}^{-1}\colon e_0(Ar)=xr,
\]
with the analogous formula at the target.
Here $\alpha_{r,A,e_0}\colon(e_0A)r=e_0(Ar)$ is the associator.
\end{lemma}
\begin{proof}[Proof idea]
Restrict the underlying comparison to $\eta\star r$.
Whisker each endpoint equation by $r$ and use compatibility of
whiskering with vertical composition and the associator. The inverse
associators in the displayed frames convert the result to exactly the
endpoint equations for the restricted expressions.
\end{proof}

\begin{lemma}
\label[lemma]{lem:Supplementary_Expression_Restriction_Laws}
Restricting an expression by $\mathrm{id}_{\Gamma}$ and retargeting its
endpoints by the right unitors gives an endpoint-preserving comparison
to the original expression. For $r\colon\Delta\to\Gamma$ and
$s\colon\Theta\to\Delta$, successive restriction by $r$ and $s$,
retargeted by $\alpha_{s,r,x}$ and $\alpha_{s,r,y}$, compares with
restriction by $r\circ s$.
\end{lemma}
\begin{proof}[Proof idea]
Use the right unitor of $A$ for the first underlying arrow comparison,
and $\alpha_{s,r,A}$ for the second. Their endpoint equations are the
unit and successive-substitution equations for transport through a
fixed functor $e_0$ or $e_1$. The finite unitor and associator compatibility calculations, derived
from the pentagon and triangle, establish these equations while retaining
the given frames.
These are the same product-calculus comparisons discussed in
\href{supplementary-pairing-and-precomposition.html}{Pairing and precomposition}.
\end{proof}

\begin{lemma}
\label[lemma]{lem:Supplementary_Composition_Substitution}
Assume the Segal axiom. For composable expressions $f\colon x\to y$,
$g\colon y\to z$, and $r\colon\Delta\to\Gamma$, there is an
endpoint-preserving comparison
\[
 (g|_r)\circ(f|_r)\simeq(g\circ f)|_r.
\]
If $a\colon x=x'$, $b\colon y=y'$, and $c\colon z=z'$ change the
endpoints, composing the retargeted expressions compares with retargeting
$g\circ f$ by $a,c$.
\end{lemma}
\begin{proof}[Proof idea]
A presentation of $g\circ f$ consists of its chosen triangle, comparisons
of the short edges with $f,g$, and the long edge with the composite.
It includes compatibility at all three vertices.
Restrict this entire presentation by $r$. The endpoint calculation of
\Cref{lem:Supplementary_Expression_Restriction} supplies its new boundary
comparisons. It is now a presentation of $(g\circ f)|_r$ as a composite
of $f|_r,g|_r$. Recognition of composites through the Segal equivalence
compares it with the chosen composite of those inputs.

For retargeting, change the three vertex frames and retain the triangle.
The common change $b$ at the middle vertex occurs on both short edges;
its compatibility is retained in the new presentation. Apply the same
recognition argument. The comparison preserves the two outer endpoints.
\end{proof}

These lemmas provide the individual substitution comparisons used in
specializing the unit and associativity arguments. They do not assert
an additional coherence theorem for all choices of Segal inverses.
\end{document}
